\documentclass[11pt,a4paper]{article}

% ---------------------------------------------------------------- packages
\usepackage[utf8]{inputenc}
\usepackage[T1]{fontenc}
\usepackage{lmodern}
\usepackage[a4paper,margin=2.2cm]{geometry}
\usepackage{xcolor}
\usepackage{enumitem}
\usepackage{amsmath,amssymb}
\usepackage{booktabs}
\usepackage{array}
\usepackage{pifont}
\usepackage{parskip}
\usepackage{hyperref}

% ------------------------------------------------------------ ownership tags
\definecolor{eddieblue}{HTML}{1F6FB2}
\definecolor{jakeorange}{HTML}{D2691E}
\definecolor{eithergrey}{HTML}{555555}

\newcommand{\eddie}{\textcolor{eddieblue}{\textsf{[E]}}}
\newcommand{\jake}{\textcolor{jakeorange}{\textsf{[J]}}}
\newcommand{\either}{\textcolor{eithergrey}{\textsf{[?]}}}

% -------------------------------------------------------------- task markers
\newcommand{\todo}{\ding{109}\,}    % open square
\newcommand{\wip}{\ding{72}\,}      % star
\newcommand{\done}{\ding{51}\,}     % check

\hypersetup{
  colorlinks=true,
  linkcolor=eddieblue,
  urlcolor=eddieblue,
  pdftitle={SwegHammer Project Handbook},
  pdfauthor={Eddie Handford and Jake},
}

% ------------------------------------------------------------------ title
\title{\textbf{SwegHammer}\\\large Project Handbook \& Living To-Do List}
\author{Eddie Handford \& Jake}
\date{Last updated: \today}

\begin{document}
\maketitle

\tableofcontents
\bigskip\hrule\bigskip

% =========================================================================
\section{How to use this document}

This is the single living reference for SwegHammer. Open it when you sit
down, close it when you walk away. It contains the project vision, the
math behind each combat phase, a checklist of everything left to do, and
ownership tags so we both know who is on what.

\paragraph{For Claude sessions:} the operational rules
(test-before-push, doc-update workflow, branch naming, git identity,
catalogue layering) live in \texttt{CLAUDE.md} at the repo root. Read
that file as well \textemdash{} it is auto-loaded into the Claude
context for every session.

\subsection*{Conventions}
\begin{itemize}[itemsep=2pt]
  \item \todo open task \quad \wip in progress \quad \done done
  \item \eddie\ Eddie's slice \quad \jake\ Jake's slice \quad \either\ unclaimed
\end{itemize}

When you pick a task up, swap its tag to your initial and flip the marker
to \wip. When you finish, flip to \done\ and commit the \texttt{.tex}
change alongside the code change so the doc never drifts from reality.

\subsection*{Build}
\begin{verbatim}
pdflatex PROJECT.tex && pdflatex PROJECT.tex
\end{verbatim}
Two passes so the table of contents resolves. No external assets needed.

% =========================================================================
\section{Project vision}

SwegHammer is a data-driven Warhammer 40K simulator with two goals:

\begin{enumerate}[itemsep=2pt]
  \item Derive fair, transparent points costs from Lanchester's Square
    Law plus empirical simulation.
  \item Test minimal rule tweaks (unit-by-unit activation,
    command-point compensation for the smaller army) that fix structural
    feel-bad moments like first-player advantage.
\end{enumerate}

The simulator is the experimental apparatus. It must be \emph{fast}
enough to run thousands of battles in a calibration sweep, and
\emph{detailed} enough to render a single battle as a watchable replay
for sanity-checking the model.

% =========================================================================
\section{Where the project stands today}

\subsection{What works}
\begin{itemize}[itemsep=2pt]
  \item \done Stochastic combat engine (\texttt{code/simulator.py}):
    roll-to-hit, AP-adjusted armour save, cover modifier, damage.
  \item \done Unit catalogue built from BSData WH40k 10e
    (\texttt{v10.6.0}) via \texttt{code/bsdata/}, with $\sim$1294 units
    across all major factions. Override layer in
    \texttt{data/overrides.json}; calibrated point costs cached in
    \texttt{data/calibrated\_points.json}.
  \item \done Unit-by-unit alternating activation
    (\texttt{Battle.\_run\_round}).
  \item \done Command-point bonus for the smaller army from Round 2.
  \item \done Streamlit dashboard (\texttt{app.py}) with preset battles,
    attrition curves, survivor histograms, points-vs-winrate sweep.
\end{itemize}

\subsection{What's missing in one sentence}

The current ``combat'' is a single abstract attack action. There are no
\emph{phases} (movement / shooting / charge / melee / morale), no map,
no positions, and no terrain beyond a binary cover flag. The rest of
this document is the plan to fill that in.

% =========================================================================
\section{Game phases}
\label{sec:phases}

40K splits a turn into discrete phases. We will mirror that structure
but keep each phase pluggable so we can disable expensive ones for fast
calibration runs.

\subsection{Command phase}
At the start of a unit's activation: gain 1 CP, resolve start-of-turn
abilities, apply the smaller-army CP compensation rule.

\paragraph{To do}
\begin{itemize}[label={},leftmargin=*,itemsep=2pt]
  \item \todo \eddie\ Add a \texttt{CommandPhase} step run at activation
    start.
  \item \todo \eddie\ Move existing CP bonus logic out of
    \texttt{Battle.\_run\_round} into the new phase so it lives in one
    place.
\end{itemize}

\subsection{Movement phase}
A unit moves up to its Move characteristic in inches. Optionally Advance
($+d6$, no shooting that turn) or Fall Back (no shooting, no charging).

\paragraph{Math.} Position update for a unit at $\mathbf{p}_t$ choosing
destination $\mathbf{p}_{t+1}$ subject to
\[
  \lVert \mathbf{p}_{t+1} - \mathbf{p}_t \rVert
  \;\leq\; M_{\text{unit}} \cdot \sigma_{\text{terrain}},
\]
where $\sigma_{\text{terrain}} \in (0,1]$ is the terrain-slowdown factor
on the line the unit traverses.

\paragraph{To do}
\begin{itemize}[label={},leftmargin=*,itemsep=2pt]
  \item \done \eddie\ Add \texttt{position: (float, float)} to
    \texttt{Unit}.
  \item \done \eddie\ Add \texttt{move: float} (inches) to
    \texttt{UnitProfile}.
  \item \done \eddie\ Implement \texttt{MovementPhase} -- move-then-shoot
    sub-phase pair per activation. Targeting is now objective-aware via
    \texttt{code/strategy.py} (HOLD / CAPTURE / STEAL / ENGAGE /
    REPOSITION intents); units sitting on objectives stay put rather
    than abandoning VP to chase the nearest enemy.
  \item \done \either\ Advance roll ($+d6$, no shooting that turn)
    implemented; range threshold is weapon range when engaging and
    objective control radius when capturing. Fall Back still
    unimplemented.
  \item \todo \either\ Segment-vs-impassable-terrain collision (only
    the destination is currently checked, so units can tunnel through
    walls).
\end{itemize}

\subsection{Shooting phase}
For every weapon in range and with line of sight: roll to hit, roll to
wound, target rolls armour save, apply damage. \textbf{The current sim
skips the wound roll}, which is the biggest single bit of math debt.

\paragraph{Math.} Expected damage of one shot:
\[
  \mathbb{E}[D] \;=\; P_{\text{hit}}
  \cdot P_w(S, T)
  \cdot \bigl(1 - P_{\text{save}}(\mathrm{sv}, \mathrm{AP}, \mathrm{cover})\bigr)
  \cdot d,
\]
where the wound probability $P_w$ follows the strength-vs-toughness
table:
\[
  P_w(S, T) =
  \begin{cases}
    5/6 & S \geq 2T \\
    4/6 & S > T \\
    3/6 & S = T \\
    2/6 & S < T \\
    1/6 & 2S \leq T
  \end{cases}
\]
and the save probability uses an effective save characteristic
$\mathrm{sv}^* = \mathrm{sv} - \mathrm{AP} - \mathrm{cover}$ capped at
$2$:
\[
  P_{\text{save}}(\mathrm{sv}^*) \;=\;
  \max\!\Bigl(0,\; \tfrac{7 - \mathrm{sv}^*}{6}\Bigr).
\]

\paragraph{To do}
\begin{itemize}[label={},leftmargin=*,itemsep=2pt]
  \item \done \eddie\ Add \texttt{strength: int} to \texttt{UnitProfile}.
  \item \done \eddie\ Add \texttt{toughness: int} to \texttt{UnitProfile}.
  \item \done \eddie\ Implement the wound table above as
    \texttt{wound\_probability(S, T)} and wire it into
    \texttt{Unit.attack()} between the hit and save rolls.
  \item \todo \eddie\ Regenerate \texttt{BASELINE.md} catalogue table to
    reflect the wound-roll-aware point costs (existing table is now
    stale; mirror match holds on symmetric terrain, but cross-faction
    matchups need recalibrating).
  \item \done \eddie\ Add \texttt{range\_inches: int} to the weapon
    profile; gate shooting on range and LoS. Liang-Barsky
    segment-vs-rectangle intersection handles obscuring terrain.
\end{itemize}

\subsection{Charge phase}
Declare a charge against an enemy within $12''$. Roll $2d6$; if the
total is at least the distance to the target, move into base contact.

\paragraph{Math.}
\[
  P_{\text{charge}}(d) \;=\; \Pr[2d_6 \geq d]
  \;=\; \frac{\#\{(a,b) \in [1,6]^2 : a+b \geq d\}}{36}.
\]
Useful values: $P(6'') = 26/36 \approx 72\%$,
$P(9'') = 10/36 \approx 28\%$, $P(12'') = 1/36 \approx 3\%$.

\paragraph{To do}
\begin{itemize}[label={},leftmargin=*,itemsep=2pt]
  \item \done \either\ Implement \texttt{ChargePhase} with declaration,
    $2d6$ roll and move-into-engagement ($1''$ of target). Lives in
    \texttt{Battle.\_do\_charge}; gated by a charge-desire heuristic
    so pure shooters don't suicide-charge out of range.
  \item \done \either\ \texttt{\_charging\_this\_round} set tracked on
    \texttt{Battle}; chargers fight first in the melee phase.
\end{itemize}

\subsection{Fight (melee) phase}
Units in base contact fight. Each model makes \texttt{attacks} attacks
at \texttt{weapon\_skill}, then wound, save, damage as per shooting.

\paragraph{To do}
\begin{itemize}[label={},leftmargin=*,itemsep=2pt]
  \item \done \either\ \texttt{UnitProfile} now carries
    \texttt{melee\_attacks}, \texttt{melee\_damage\_per\_shot},
    \texttt{melee\_hit\_probability}, \texttt{melee\_strength},
    \texttt{melee\_ap}, \texttt{melee\_weapon}. Mapper extracts
    \texttt{typeName="Melee Weapons"} profiles and picks a best-legal
    loadout; 1261 of 1291 units carry a usable melee profile.
  \item \done \either\ \texttt{MeleePhase} (\texttt{Battle.\_do\_fight}):
    any unit within $1.5''$ of an enemy strikes with its melee profile;
    chargers fight first. \texttt{Unit.attack(\dots,\, mode="melee")}
    switches the stat block.
\end{itemize}

\subsection{Morale / battleshock phase}
A unit that lost models tests Leadership; a failed test inflicts further
casualties or status effects.

\paragraph{To do}
\begin{itemize}[label={},leftmargin=*,itemsep=2pt]
  \item \done \either\ \texttt{UnitProfile} carries
    \texttt{leadership} (Ld) and \texttt{oc} (Objective Control); both
    are extracted by the mapper from the BSData unit profile.
  \item \done \either\ 10e Battleshock implemented in the simulator:
    from Round 2, any unit below half HP rolls $2d6$ vs Ld; a fail
    leaves it Battleshocked for the round (OC drops to 0 for objective
    scoring). Detachment \texttt{enemy\_ld\_penalty} composes into the
    target.
\end{itemize}

% =========================================================================
\section{Terrain and map}
\label{sec:terrain}

\subsection{Coordinate system}
Continuous 2D plane, units in inches. Standard boards: $44'' \times
60''$ (Combat Patrol), $44'' \times 90''$ (Strike Force). Models are
circles with a defined base radius.

\subsection{Terrain types}
\begin{tabular}{lll}
\toprule
Type & Effect & Notes \\
\midrule
Open        & none                              & default \\
Light cover & $+1$ save vs ranged               & matches today's ``in cover'' flag \\
Heavy cover & $+1$ save and $-1$ to hit         & ruins, barricades \\
Obscuring   & blocks LoS through it             & forests, walls \\
Impassable  & blocks movement                   & cliffs \\
\bottomrule
\end{tabular}

\subsection{Line of sight}
Naive implementation: draw a segment from the attacker centre to the
target centre, check intersection with any obscuring polygon. If that
becomes a perf bottleneck, swap in a spatial hash over terrain
polygons.

\paragraph{To do}
\begin{itemize}[label={},leftmargin=*,itemsep=2pt]
  \item \done \eddie\ \texttt{Terrain} dataclass (axis-aligned
    rectangles for now; polygon upgrade can wait).
  \item \done \eddie\ \texttt{Map} dataclass: board size, terrain list,
    deployment-zone depth, and objective markers. Every stock map
    ships a 5-objective quincunx (centre $+$ four equidistant markers
    $\sim$30\% in from each edge), matching the geometry of 10e Take
    and Hold / Crucible / Tipping Point. \texttt{Objective} dataclass
    in \texttt{code/map.py}: \texttt{(x, y, control\_radius=3'',
    vp\_per\_round=5)}.
  \item \done \eddie\ \texttt{has\_line\_of\_sight(attacker, target)}
    on \texttt{Map}, Liang-Barsky parametric clipping over obscuring
    rectangles. Units inside obscuring terrain can still shoot out and
    be shot at (the containing rectangle is excluded from the
    intersection test).
  \item \done \eddie\ \texttt{cover\_at(point)} on \texttt{Map}
    returning the strongest cover the point sits inside; terrain-aware
    cover is applied per-shot in \texttt{Battle.\_do\_shoot}.
  \item \done \eddie\ Three stock maps registered in
    \texttt{STOCK\_MAPS}: \texttt{combat\_patrol},
    \texttt{open\_plains}, \texttt{urban\_sprawl}. The first two are
    point-symmetric reference maps; Urban Sprawl is deliberately
    asymmetric.
\end{itemize}

% =========================================================================
\section{Visualization: the dual-mode plan}
\label{sec:viz}

The simulator needs to do two things that pull in opposite directions:
\begin{enumerate}[itemsep=2pt]
  \item Run thousands of fast battles for calibration sweeps.
  \item Render a single battle in detail for sanity checks and demos.
\end{enumerate}

\paragraph{Design.} Keep the simulation core
(\texttt{simulator.py}, \texttt{army.py}, \texttt{units.py})
\emph{completely headless and deterministic given a seed}. Build
rendering on top as an event-subscriber layer:

\begin{itemize}[itemsep=4pt]
  \item \textbf{Detailed mode} -- a \emph{Pygame} or
    \emph{matplotlib-animation} renderer with units drawn as tokens on
    a 2D top-down map, with range circles, LoS lines, charge arcs. Used
    for: replaying a single battle, demo videos, debugging movement bugs.
  \item \textbf{Fast mode} -- \emph{no rendering at all}, just numerical
    state. Used for: Monte Carlo sweeps where we run $10^3$--$10^4$
    battles per matchup and only care about win rates and attrition
    curves.
\end{itemize}

Crucially, the two modes share \emph{the same} \texttt{Battle} object;
the mode flag just toggles which subscriber is attached. A bug in
the visualisation therefore cannot affect calibration numbers.

\paragraph{To do}
\begin{itemize}[label={},leftmargin=*,itemsep=2pt]
  \item \done \eddie\ Refactor \texttt{Battle.run()} to emit events
    via a subscriber list. Events defined in \texttt{code/events.py}:
    \texttt{BattleStarted}, \texttt{RoundStarted}, \texttt{UnitActivated},
    \texttt{UnitMoved}, \texttt{UnitShot}, \texttt{UnitKilled},
    \texttt{RoundEnded}, \texttt{BattleEnded}.
  \item \done \eddie\ Null-subscriber fast mode -- empty subscriber list
    means zero per-event overhead in calibration sweeps.
  \item \todo \either\ \texttt{PygameRenderer} -- consumes the event
    stream, paints frames live.
  \item \done \eddie\ Matplotlib static-frame renderer
    (\texttt{code/renderer.py}) reconstructing world state from the
    event stream; feeds the Streamlit replay slider page. Renders
    advance / charge / melee / objective events with directional
    arrows, star bursts on hits, and gold rings around scoring
    objective markers.
\end{itemize}

% =========================================================================
\section{Front end}

The current Streamlit app (\texttt{app.py}) is a good v1 dashboard. We
will extend it rather than replace it.

\paragraph{Near-term}
\begin{itemize}[label={},leftmargin=*,itemsep=2pt]
  \item \done \eddie\ Embed a scrub-through replay of one battle on a
    map (Streamlit "Watch a battle" tab; matplotlib top-down view with
    deployment-zone hints, terrain tiles, unit dots, move/shot arrows
    on the current event).
  \item \done \eddie\ Map picker -- sidebar dropdown over
    \texttt{STOCK\_MAPS} (one map for now; structure ready for more).
  \item \todo \eddie\ Per-unit cover and terrain overrides in the
    sidebar (current cover flag is army-wide).
\end{itemize}

\paragraph{Longer-term}
\begin{itemize}[label={},leftmargin=*,itemsep=2pt]
  \item \todo \either\ Army builder UI: drag-and-drop datasheets, points
    totalled live.
  \item \todo \either\ Calibration browser: view the most recent sweep,
    drill into outlier matchups.
\end{itemize}

% =========================================================================
\section{Datasheet ingestion}

Importing real datasheets at scale is Jake's slice of the project. The
shape of the data model needs to be agreed before he starts, otherwise
we will end up rewriting \texttt{units.py} twice.

\paragraph{Status: v1 complete (BSData WH40k 10th-ed, $\sim$1100 units).}

\begin{itemize}[label={},leftmargin=*,itemsep=2pt]
  \item \done \jake\ Pick a source format \textemdash{} BSData WH40k 10e
    BattleScribe XML files, pinned to release \texttt{v10.6.0}. See
    \texttt{code/bsdata/fetch.py} for the GitHub-backed cache layer.
  \item \done \jake\ Define a target schema \textemdash{}
    \texttt{MappedUnit} (\texttt{code/bsdata/mapper.py}) lowers to
    \texttt{UnitProfile} via the loader in
    \texttt{code/bsdata/loader.py}, merging \texttt{parsed.json} with
    user-tunable \texttt{data/overrides.json}.
  \item \done \jake\ Build a loader yielding \texttt{UnitProfile}
    instances \textemdash{} \texttt{UNIT\_CATALOG} is now built at
    import time from the merged BSData + overrides view.
  \item \done \jake\ Multi-weapon flattening \textemdash{} the mapper
    picks the single best-legal weapon by
    $A \cdot P_{\text{hit}} \cdot D \cdot (1 - P_{\text{save\_baseline}})$,
    i.e.\ the loadout maximising expected damage through baseline-Marine
    armour. List of alternatives is kept in
    \texttt{MappedUnit.loadout} for audit.
\end{itemize}

\paragraph{Phase 2 follow-ups (see \texttt{ROADMAP.md}).}
\begin{itemize}[label={},leftmargin=*,itemsep=2pt]
  \item \todo \jake\ Extract \texttt{S} (weapon Strength) and \texttt{T}
    (unit Toughness) from BSData into \texttt{UnitProfile} \textemdash{}
    currently defaulted to S$=$T$=4$ via \texttt{BASELINE\_*}.
  \item \todo \jake\ Fall-back to melee weapons for units with no ranged
    profile ($\sim$260 currently flagged \texttt{enabled: false}).
  \item \todo \jake\ Decide per-squad vs per-model aggregation
    for \texttt{health} (currently per-model wounds).
\end{itemize}

% =========================================================================
\section{Combat math reference card}

\subsection{Hit / wound / save chain}
For one attack with weapon profile $(P_{\text{hit}}, S, \mathrm{AP}, d)$
against target $(T, \mathrm{sv})$ and cover bonus $c \in \{0,1\}$:
\[
  \mathbb{E}[\text{damage}]
  \;=\; P_{\text{hit}}
  \cdot P_w(S, T)
  \cdot \bigl(1 - P_{\text{save}}(\mathrm{sv} - \mathrm{AP} - c)\bigr)
  \cdot d.
\]

\subsection{Lanchester score}
\[
  \text{Score}(\text{unit}) \;=\; \bigl(H \cdot D \cdot P_{\text{hit}}\bigr)^2,
  \qquad
  \text{Score}(\text{army}) \;=\; \sum_i e_i^{\,2}.
\]

\subsection{Charge probability}
\[
  P_{\text{charge}}(d)
  \;=\; \frac{1}{36}
    \sum_{a=1}^{6} \sum_{b=1}^{6} \mathbf{1}\!\left[a+b \geq d\right].
\]

\subsection{Primary VP scoring and win condition}
At end of each round, for every objective $o$ on the map let
$\mathrm{OC}^A_o, \mathrm{OC}^B_o$ be the total Objective Control of
each side's alive units inside the objective's control radius
(Battleshocked units contribute $\mathrm{OC} = 0$). The side with the
strict majority banks $o.\text{vp\_per\_round}$ Primary VP that round.
\[
  \mathrm{VP}^A \mathrel{+}= \sum_{o \in O}
    o.v \cdot \mathbf{1}\!\left[\mathrm{OC}^A_o > \mathrm{OC}^B_o\right],
  \qquad
  \mathrm{VP}^B \text{ symmetrically}.
\]
Win condition (replaces ``most surviving units''):
\begin{enumerate}[itemsep=2pt]
  \item One-sided wipe: if only one side has alive units, that side wins.
  \item Else: higher Primary VP wins.
  \item VP tie: higher remaining army points wins ($10\%$ margin else
    draw).
\end{enumerate}

\subsection{Strategy layer}
Movement targeting is no longer ``march at the nearest enemy''. The
strategy layer (\texttt{code/strategy.py}) selects a $(\text{target\_pos},
\text{intent})$ per activation from five intents: \textsc{Hold},
\textsc{Capture}, \textsc{Steal}, \textsc{Engage}, \textsc{Reposition}.
Objectives are scored by current control state (enemy-held $\times 3.5$
to bias steals, uncontested $\times 2.5$, ours $\times 1.0$) and
distance-weighted; role bias (\texttt{code/roles.py}, classifying units
into \textsc{Shooty} / \textsc{Melee} / \textsc{Dual} / \textsc{Horde}
/ \textsc{Heavy} / \textsc{Support}) tips ties towards holding when in
range vs always-closing.

% =========================================================================
\section{Collaboration workflow}

\subsection{Ownership at a glance}
\begin{tabular}{lll}
\toprule
Area & Lead & Status \\
\midrule
Simulation core (phases, math)    & \eddie\ Eddie & in progress \\
Datasheet ingestion + catalogue   & \jake\ Jake   & v1 complete \\
Terrain \& maps                   & \either\ TBD  & not started \\
Visualisation (dual-mode)         & \either\ TBD  & not started \\
Streamlit front end               & \eddie\ Eddie & v1 done \\
\bottomrule
\end{tabular}

\subsection{Picking the project back up}
\begin{enumerate}[itemsep=2pt]
  \item \texttt{git pull}.
  \item Open this PDF and scan the unchecked items in your area.
  \item Pick the smallest task that unblocks something else.
  \item Commit with a clear message; tick the box in this file in the
    same commit.
\end{enumerate}

\subsection{When something is safe to merge to main}
\begin{itemize}[itemsep=2pt]
  \item \texttt{python -m unittest discover tests} must pass (currently
    73 tests, growing).
  \item \texttt{python run.py --cli} must exit 0 (demo plus 100-battle
    calibration complete cleanly).
  \item \textbf{Mirror Intercessor matchup VP totals:} across 200 battles
    at 1000 pts on \texttt{combat\_patrol}, neither side's average VP
    exceeds the other by more than ${\sim}25\%$. This replaces the old
    win-rate-spread gate because the new combat model (5-round +
    objectives + abilities) produces more draws and decisive blowouts
    than the old attrition-only sim, so raw win-rate is no longer a
    sensitive symmetry probe.
  \item \textbf{No new units silently disabled:} run
    \texttt{python -m code.bsdata.audit} and confirm the skip-reason
    histogram is unchanged from the previous commit (or any new skips
    are accounted-for in the commit message).
  \item The Streamlit app still launches and the preset battles still
    produce charts and a watchable replay.
\end{itemize}

% =========================================================================
\section{Glossary}

\begin{tabular}{lp{0.7\linewidth}}
\toprule
Term & Meaning \\
\midrule
AP                & Armour penetration; negative integer degrading target save. \\
BS / WS           & Ballistic / Weapon Skill; roll needed to hit. \\
Charge phase      & Phase where units roll $2d6$ to close into base contact. \\
Cover             & Terrain bonus to saving throws or hit penalty. \\
CP                & Command Point; resource for stratagems and rerolls. \\
Lanchester score  & $(H \cdot D \cdot P)^2$; theoretical effectiveness of a unit. \\
LoS               & Line of sight; whether the attacker can ``see'' the target through terrain. \\
Stochastic mode   & Simulator with real dice rolls. \\
Deterministic mode& Simulator using expected values, no rolls. \\
\bottomrule
\end{tabular}

\end{document}
